\section{Welcome Letter from the Chairs}

\begin{chairletter}[Chairs of the United Nations Security Council IEUMUN 2026]{Mahima Rajpal, Talia Asar, and Aryan Agrawal}
Dear Distinguished Delegates,

Welcome to the United Nations Security Council at IEUMUN 2026. We are Mahima Rajpal, Talia Asar and Aryan Agrawal, and we are delighted to welcome you to one of the conference's most dynamic and intellectually demanding committees. First, we would like to introduce ourselves.

I’m Mahima Rajpal, an Electrical Engineering and IT Major at the Karlsruhe Institute of Technology (KIT), Germany. Although my academic background is rooted in engineering, I have always been fascinated by international affairs and the role diplomacy plays in addressing complex global challenges. MUNs have become the perfect intersection of these interests, combining analytical thinking with negotiation, leadership, and public speaking. Having participated in and led conferences across the European MUN circuit in various roles, I am excited to chair this year's Security Council alongside my co-chairs. What I enjoy most about the Security Council is its ability to bring together legal, political, economic, and strategic perspectives to address some of the world's most pressing challenges. Maritime security embodies exactly that complexity, and I hope this committee encourages you to think beyond immediate crises and develop practical, balanced, and lasting solutions.

My name is Ashanti Talia Asar and I have just graduated with my bachelors in Business Administration from Hult International Business school. I first started MUN in high school and enjoyed it so much that I stuck to it. Since then I have had the privilege of joining HultMUN first as a delegate, then leading for a year as a president and now being a proud HultMUN alumni. I have always loved security related and financial committees. Last year I chaired the WTO at IEUMUN and this year I have the pleasure of chairing UNSC. I am very excited to see what all of you will bring to the debate and what solutions you will bring to the table. I hope this committee pushes you to think outside the box in ways that encourages interdisciplinary thinking.

I am Aryan Agrawal, a second-year student pursuing a Dual Degree in Economics and International Relations at IE University in Spain. My journey with debate and Model United Nations began at the age of nine, and since 2018, I have had the privilege of actively delegating, chairing internationally, and organising conferences across different platforms. Over the years, I have come to deeply value MUN for the way it creates a space where people can open up, exchange ideas, and engage meaningfully with one another. For me, MUN is not just about procedure or debate, but about the unique blend of diplomacy, critical thinking, collaboration, and leadership that it fosters. These are qualities that continue to inspire me and shape the way I approach every committee I am part of. I am truly honoured to serve as a chair at IEUMUN 2026, and I look forward to working with all delegates in a dynamic, engaging, and productive committee.

This Study Guide provides the foundation for your preparation, but it should not replace independent research. We encourage you to examine your country’s strategic interests, trade and energy dependencies, interpretation of international maritime law, relationships with regional actors, and position on naval or multilateral security mechanisms.

A clear and well-researched Position Paper will help you enter a committee with realistic priorities, defensible red lines, and workable proposals. The strongest delegates will not merely describe the problem; they will understand what their state can support, what it will resist, and where compromise remains possible.

We are excited to welcome you to IEUMUN 2026. Come prepared to question assumptions, defend your mandate, negotiate strategically, and confront one of the most consequential security challenges currently being faced by the world.

\signoff[Warm regards,]{Mahima Rajpal, Talia Asar, and Aryan Agrawal}
\end{chairletter}
